% ============================================================
\section{Threat Model}

\subsection{Attacker Capabilities}

We assume a controlled research setting in which the attacker can inject malicious or misleading content into one or more untrusted channels, such as:

\begin{itemize}
    \item Retrieved documents.
    \item Tool outputs.
    \item Emails or task inputs.
    \item Inter-agent messages.
    \item External knowledge sources.
\end{itemize}

The attacker does not access real private data. All sensitive information used in the experiment is synthetic and represented using canary secrets.

\subsection{Attack Goals}

We consider two attack goals.

\subsubsection{Sensitive Information Leakage}

The attacker attempts to cause the system to expose a synthetic secret, either in the final output or in an internal system channel.

\subsubsection{Performance Degradation}

The attacker attempts to reduce task performance or increase operational cost without necessarily leaking a secret.

\subsection{Out of Scope}

This work does not target real users, real private data, production systems, credential theft, or unauthorized access. All experiments are performed in sandboxed environments using synthetic secrets and controlled tasks.

% ============================================================
\section{Experimental Systems}

We evaluate the framework on four systems.

\subsection{AgentDojo}

AgentDojo serves as the benchmark-oriented baseline environment. It provides realistic tool-using agent tasks and is suitable for evaluating prompt-injection robustness, utility preservation, and attack success.

\subsection{AutoGen}

AutoGen represents a conversational multi-agent framework in which agents interact through messages and may use tools or human-in-the-loop workflows.

\subsection{CrewAI}

CrewAI represents role-based multi-agent collaboration, where agents are assigned goals, roles, and tools, and may delegate subtasks to one another.

\subsection{LangGraph}

LangGraph represents stateful graph-based agent workflows. It is especially useful for tracing where an attack succeeds, because intermediate nodes, edges, state updates, and tool calls can be logged explicitly.

\subsection{Unified Experimental Scenario}

To enable comparison, each framework is configured to implement a similar abstract workflow:

\begin{figure}[h]
\centering
\fbox{
\begin{minipage}{0.9\linewidth}
\centering
User Task $\rightarrow$ Planner Agent $\rightarrow$ Retriever/Tool Agent $\rightarrow$ Worker Agent $\rightarrow$ Reviewer/Guard Agent $\rightarrow$ Final Output
\end{minipage}
}
\caption{Unified multi-agent workflow used across frameworks.}
\label{fig:workflow}
\end{figure}

% ============================================================
\section{AutoResearch Red-Teaming Framework}

\subsection{Overview}

The proposed framework uses AutoResearch-style optimization to generate, test, and select attack variants. Each loop receives a target system, a task, an attack space, and a scoring function.

\begin{figure}[h]
\centering
\fbox{
\begin{minipage}{0.9\linewidth}
\centering
Generate Attack Variant $\rightarrow$ Run Evaluation $\rightarrow$ Score Result $\rightarrow$ Keep or Reject $\rightarrow$ Next Iteration
\end{minipage}
}
\caption{AutoResearch red-teaming loop.}
\label{fig:autoresearch_loop}
\end{figure}

\subsection{Loop A: Leakage Optimization}

The first loop optimizes for sensitive-information leakage.

\begin{equation}
\text{Objective}_{leak} = \max \left( \text{TotalExposureRate} \right)
\end{equation}

The loop may modify:

\begin{itemize}
    \item Attack prompt wording.
    \item Injection location.
    \item Target agent.
    \item Tool-output content.
    \item Retrieved document content.
    \item Canary placement.
    \item Attack stealth level.
\end{itemize}

\subsection{Loop B: Performance Degradation Optimization}

The second loop optimizes for reduced utility or increased cost.

\begin{equation}
\text{Objective}_{perf} = \max \left( \alpha \cdot \text{UtilityDrop} + \beta \cdot \text{CostAmplification} \right)
\end{equation}

where $\alpha$ and $\beta$ control the relative importance of correctness degradation and cost increase.

The loop may modify:

\begin{itemize}
    \item Misleading documents.
    \item Ambiguous task inputs.
    \item Tool outputs.
    \item Delegation-inducing instructions.
    \item Retry-inducing perturbations.
    \item Conflicting context.
\end{itemize}

% ============================================================
\section{Evaluation Metrics}

\subsection{Leakage Metrics}

\subsubsection{Leakage Attack Success Rate}

\begin{equation}
\text{Leakage ASR} =
\frac{\#\text{runs with leakage}}{\#\text{attacked runs}}
\end{equation}

\subsubsection{Final-Output Leakage Rate}

\begin{equation}
\text{FinalOutputLeakage} =
\frac{\#\text{runs with leakage in final output}}{\#\text{attacked runs}}
\end{equation}

\subsubsection{Internal Leakage Rate}

\begin{equation}
\text{InternalLeakage} =
\frac{\#\text{runs with leakage in internal channels}}{\#\text{attacked runs}}
\end{equation}

\subsubsection{Total Exposure Rate}

\begin{equation}
\text{TotalExposureRate} =
\frac{\#\text{runs with leakage in any channel}}{\#\text{attacked runs}}
\end{equation}

\subsubsection{Output-Only Miss Rate}

\begin{equation}
\text{OutputOnlyMissRate} =
\frac{\#\text{runs with internal leakage but no final-output leakage}}{\#\text{attacked runs}}
\end{equation}

\subsection{Performance Metrics}

\subsubsection{Utility Drop}

\begin{equation}
\text{UtilityDrop} =
\text{CleanTaskSuccessRate} - \text{AttackedTaskSuccessRate}
\end{equation}

\subsubsection{Cost Amplification}

\begin{equation}
\text{CostAmplification} =
\frac{\text{AttackedCost}}{\text{CleanCost}}
\end{equation}

\subsubsection{Latency Increase}

\begin{equation}
\text{LatencyIncrease} =
\frac{\text{AttackedLatency}}{\text{CleanLatency}}
\end{equation}

\subsubsection{Tool Call Increase}

\begin{equation}
\text{ToolCallIncrease} =
\frac{\text{AttackedToolCalls}}{\text{CleanToolCalls}}
\end{equation}

\subsection{Transferability Metrics}

\subsubsection{Transfer ASR}

\begin{equation}
\text{TransferASR} =
\frac{\text{ASR on unseen systems}}{\text{ASR on source system}}
\end{equation}

\subsubsection{Generalization Gap}

\begin{equation}
\text{GeneralizationGap} =
\text{ASR}_{seen} - \text{ASR}_{unseen}
\end{equation}

\subsection{Defense Metrics}

\begin{itemize}
    \item Leakage reduction.
    \item Utility preservation.
    \item False refusal rate.
    \item Over-defense rate.
    \item Defense cost overhead.
    \item Attack success despite guard.
\end{itemize}

% ============================================================
\section{Experimental Design}

\subsection{Baselines}

We compare four attack settings:

\begin{itemize}
    \item \textbf{B0: Clean}: no attack.
    \item \textbf{B1: Random}: random or naive injection.
    \item \textbf{B2: Manual}: manually designed attack.
    \item \textbf{B3: AutoResearch}: automatically optimized attack.
\end{itemize}

\subsection{Defense Settings}

We evaluate the following defense configurations:

\begin{itemize}
    \item \textbf{D0: No defense}.
    \item \textbf{D1: Prompt-based defense}.
    \item \textbf{D2: Guard agent}.
    \item \textbf{D3: Memory or inter-agent redaction}.
    \item \textbf{D4: Tool-output filtering}.
\end{itemize}

\subsection{Experimental Matrix}

Each experiment is defined by:

\begin{equation}
E = (S, G, A, D, T)
\end{equation}

where $S$ is the system, $G$ is the attack goal, $A$ is the attack type, $D$ is the defense setting, and $T$ is the task/domain.

\begin{table*}[t]
\centering
\caption{Experimental matrix.}
\label{tab:exp_matrix}
\begingroup
\small
\setlength{\tabcolsep}{3pt}
\begin{tabular}{@{}p{0.14\textwidth}p{0.18\textwidth}p{0.13\textwidth}p{0.36\textwidth}p{0.10\textwidth}@{}}
\toprule
System & Architecture & Goal & Attack Type & Defense \\
\midrule
AgentDojo & Single-agent/tool-using & Leakage & Clean / Random / Manual / AutoResearch & D0--D4 \\
AutoGen & Multi-agent & Leakage & Clean / Random / Manual / AutoResearch & D0--D4 \\
CrewAI & Multi-agent & Degradation & Clean / Random / Manual / AutoResearch & D0--D4 \\
LangGraph & Multi-agent graph & Degradation & Clean / Random / Manual / AutoResearch & D0--D4 \\
\bottomrule
\end{tabular}
\endgroup
\end{table*}
